% !Mode:: "TeX:UTF-8"

\chapter{绪论}

\section{课题背景及研究的目的和意义}

发展国防工业、微电子工业等尖端技术需要精密和超精密的仪器设备。气体静压轴承由于具有运动精度高、摩擦损耗小和无污染等特点，已在航空航天、半导体制造和精密测量领域得到应用。

本文围绕局部多孔质气体静压轴承的静态特性与稳定性开展研究，目标是建立可用于工程设计的分析模型，并通过实验数据对模型进行验证。

\section{气体润滑轴承及其相关理论的发展概况}

气体轴承是利用气膜支撑负荷或减少摩擦的机械构件。根据间隙内气膜压力的产生原理，气体轴承可以分为气体静压轴承、气体动压轴承、气体动静压轴承和气体压膜轴承等形式。

\subsection{气体润滑轴承的发展}

早期研究主要集中在气膜压力分布和承载能力分析。随着计算流体力学和精密制造技术的发展，研究重点逐渐扩展到复杂结构轴承的数值求解、稳定性分析和工程应用验证\cite{hithesis2017,Chen1992}。

\subsection{气体润滑轴承的分类}

图\ref{fig:bearing-types} 给出了气体润滑轴承分类的示意写法。图题应位于图下方，图序按章编号，图名与图序之间空两个半角字符。

\begin{figure}[htbp]
  \centering
  \fbox{\rule{0pt}{32mm}\rule{105mm}{0pt}}
  \caption{气体润滑轴承的分类}
  \label{fig:bearing-types}
\end{figure}

\subsubsection{多孔质静压轴承的分类}

多孔质静压轴承可按照供气方式、节流结构和工作面形态进行分类。四级标题可按需要使用，但本科毕业论文目录通常列到第三级标题。

\section{本文的主要研究内容}

本文主要研究内容包括：建立多孔质材料渗透率模型；建立局部多孔质气体静压止推轴承静态特性的数学模型；通过实验数据验证模型的有效性。

\chapter{基于数值方法的轴承静态特性研究}

\section{引言}

利用数值方法研究流场，可以避免对复杂偏微分方程进行完全解析求解。对于局部多孔质气体静压轴承，数值方法能够直观给出压力分布、承载能力和结构参数之间的关系。

\subsection{边界条件的设定}

本文采用连续性方程和 Darcy 定律描述多孔质区域内的流动关系。典型控制方程可写为
\begin{equation}
  Q = \frac{kA}{\mu L}\Delta P ,
  \label{eq:darcy}
\end{equation}
式中，$Q$ 为体积流量，$k$ 为渗透率，$A$ 为流通面积，$\mu$ 为动力黏度，$L$ 为流动路径长度，$\Delta P$ 为压力差。

\section{本章小结}

本章给出了局部多孔质气体静压轴承静态特性研究的数值建模思路，并说明了关键边界条件和主要物理量的含义。

\chapter{局部多孔质静压轴承的试验研究}

\section{引言}

实验研究用于验证理论模型和数值结果的合理性。本章给出多孔质石墨试样的渗透率测试示例。

\section{多孔质石墨渗透率测试试验}

1 号试样的试验数据见表\ref{tab:permeability-data}。表题应位于表上方，表序按章编号；当表格跨页时，可采用续表形式。

\begin{table}[htbp]
  \centering
  \caption{1号试样渗透率测试数据}
  \label{tab:permeability-data}
  \begin{tabular}{cccccc}
    \hline
    供气压力 & 流量测量 & 流量修正值 & 压力差 & $\lg\Delta P$ & $\lg M$ \\
    $P_s$ (MPa) & $M'$ (m$^3$/h) & $M$ ($10^{-4}$m$^3$/s) & $\Delta P$ (Pa) & & \\
    \hline
    0.15 & 0.009 & 0.02312 & 46900  & 4.67117 & -5.63601 \\
    0.20 & 0.021 & 0.04584 & 96900  & 4.98632 & -5.33876 \\
    0.25 & 0.039 & 0.07413 & 146900 & 5.16702 & -5.13001 \\
    0.30 & 0.097 & 0.16747 & 196900 & 5.29424 & -4.77606 \\
    \hline
  \end{tabular}
\end{table}

\section{本章小结}

本章给出了多孔质石墨渗透率测试的数据组织方式，并展示了本科毕业论文中表格、公式和正文引用的基本写法。
